% page 346 of the facsimile (image p-345 of the scan)
% block 154.9 x 233.3 mm ; left margin 8.0 mm
\pgc{-12.700mm}{0.000mm}{
\l{\cel{2}338.}
\l{}
\l{\sou{ligustro}. (bot.)\cel{2}Arboreto ek la familio \textquotedbl{}oleacei\textquotedbl{}, de qua}
\l{\cel{3}onu formacas masivi. - L.\cel{1}\sou{ligustrum.}DISL.}
\l{}
\l{likar. (\sou{netrans}.) (\sou{aludante generale liquido}). Eskapar tra}
\l{\cel{3}fenduro (de vazo, e c.) - DE.}
\l{}
\l{likantropo. (\sou{patol.}) Persono qua afektesas da la morbo psi-}
\l{\cel{3}kala pro qua lu imaginas transformesir a volfo. - DEFIS.}
\l{}
\l{\sou{likeno}. (\sou{bot.)}\cel{4}Planto kriptogama qua kreskas sur la kor-}
\l{\cel{3}tico di la arbori, sur la roki, sur la petri humida,}
\l{\cel{3}forme di krusto pulveratra. - L.\cel{1}\sou{cetraria}.\cel{2}DEFIS.}
\l{}
\l{\sou{likopodio.}\cel{2}(\sou{bot.}) Planto kriptogama, di qua la kapsuli kon-}
\l{\cel{3}tenas pulvero inflamebla, uzata medicinale kom sikirivo.}
\l{\cel{3}- L.\cel{1}\sou{lycopodium}.\cel{2}- DEFIS.}
\l{}
\l{\sou{liktoro}. (\sou{epoki antiqua, en Roma}). Guardano qua iris avan}
\l{\cel{3}la oficieri judiciala, portis hakilo plasizita meze di}
\l{\cel{3}vergo-fasko ed havis kom tasko exekutar la verdikto qua}
\l{\cel{3}kondamnis la krimininti a senkapigeso pos batesir per ver-}
\l{\cel{3}gi. - DEFIS.}
\l{}
\l{\sou{lilaco}. (\sou{bot.})\cel{2}Arboreto di la genero \textquotedbl{}seringo\textquotedbl{}, di qua la}
\l{\cel{3}flori dispozesas grape. - L.\cel{1}\sou{Syringa.}}
\l{}
\l{\sou{lilio}. (\sou{bot.)}\cel{2}Planto herbacea, bulboza kun stipo rekta,}
\l{\cel{3}e di qua la varietato maxim konocata havas floro pure-}
\l{\cel{3}blanka e klosho-forma. - L.\cel{1}\sou{lilium.}\cel{2}- DEFIRS.}
\l{}
\l{\sou{liliaceo.}\cel{2}(\sou{bot.)}\cel{2}Di qua la naturo esas olta di lilio}
\l{\cel{3}(planto monokotiledona). - DEFIS.}
\l{}
\l{\sou{limar}. (\sou{trans., per}).Konsumar, planigar per froto-kontakto}
\l{\cel{3}kun lamo de fero, de stalo, kun strii transversa qui as-}
\l{\cel{3}pektas quale rangi de denti. - FS.}
\l{}
\l{\sou{limako}. (\sou{zool.)}\cel{2}Molusko gasteropoda, reptera, di qua la}
\l{\cel{3}korpo esas oblonga, sen-konka, kovrata da humoro viskoza.}
\l{\cel{3}- L.\cel{1}\sou{limax}. - EF.}
\l{}
\l{\sou{limando}. (\sou{zool.}) Fisho plata, ek la genero \textquotedbl{}plii\textquotedbl{}, di qua}
\l{\cel{3}la pelo esas rugoza, aspera. - L.\cel{1}\sou{pleuronectes.}\cel{2}- FIL.}
\l{}
\l{\sou{limbo}. I. (\sou{bot.)}\cel{2}Parto supra, ordinare expansita, di la}
\l{\cel{3}koroli monopetala, di la kalici monofila. - II. (\sou{tekn.})}
\l{\cel{3}Bordo gradizita, di la cirklo uzata por mezurar la anguli,}
\l{\cel{3}- EFIRS.}
\l{}
\l{limeto. (bot.)Frukto di arboro ek la genero \textquotedbl{}oranjieri\textquotedbl{},}
\l{\cel{3}qua produktas\cel{1}\sou{mikra}\cel{1}citrono dolca. - L. citrus limeta}
\l{\cel{3}DEFISL.}
}
